\documentclass[11pt]{article}

\usepackage[margin=1in]{geometry}
\usepackage{amsmath}
\usepackage{amssymb}
\usepackage{mathtools}
\usepackage{enumitem}
\usepackage{parskip}
\usepackage{graphicx}
\usepackage{xcolor}
\usepackage{hyperref}

\newcommand{\flops}{\text{FLOPs}}
\newcommand{\flopss}{\text{FLOPs/s}}
\newcommand{\Tmath}{T_{\text{math}}}
\newcommand{\Tcomm}{T_{\text{comm}}}
\newcommand{\Tcomp}{T_{\text{comp}}}
\newcommand{\Wici}{W_{\text{ici}}}
\newcommand{\Whbm}{W_{\text{hbm}}}

\title{Scaling Book --- Section 7 Exercises}
\author{Rahul Bir}
\date{}

\begin{document}
\maketitle

\section*{Pop Quiz}

Take gen step on BS of 4 tokens on 30B param dense model on TPU v5e $4 \times 4$ slice in int8 w/ bf16 FLOPs, 8192 context, 100 kB/token KV caches.

\subsection*{(i)}

\[
T_{\text{min,step}} = \frac{4 \cdot 100 \cdot 10^3 \cdot 8192}{4 \cdot 4 \cdot 8.2 \cdot 10^{11}} + \max\!\left( \underbrace{\frac{2 \cdot 4 \cdot 30 \cdot 10^9}{4 \cdot 4 \cdot 1.97 \cdot 10^{14}}}_{7.6 \cdot 10^{-5}\text{ s}},\; \underbrace{\frac{30 \cdot 10^9}{4 \cdot 4 \cdot 8.2 \cdot 10^{11}}}_{2.2 \text{ ms}} \right)
\]
\[
\approx 2.5 \text{ ms.}
\]

MLPs are comm-bounded.

We could also see this since our total token BS is 4 per fwd pass $\ll B_{\text{crit}} = 240$.

\subsection*{(ii)}

$B = 256 > B_{\text{crit}} = 240 \Rightarrow$ MLPs are comp bound.

For thoroughness:
\[
T_{\text{min,step}} = \underbrace{\frac{256 \cdot 100 \cdot 10^3 \cdot 8096}{16 \cdot 8.2 \cdot 10^{11}}}_{\approx\, 16 \text{ ms}} + \max\!\left( \underbrace{\frac{2 \cdot 256 \cdot 30 \cdot 10^9}{16 \cdot 1.97 \cdot 10^{14}}}_{\approx\, 4.9 \text{ ms}},\; \underbrace{\frac{30 \cdot 10^9}{16 \cdot 8.2 \cdot 10^{11}}}_{\approx\, 4.6 \text{ ms}} \right)
\]
\[
= 20.9 \text{ ms}.
\]

\section*{Q1}

\subsection*{(i)}

Per layer:
\begin{align*}
\text{Attn:} \quad & 2DNH + 2DKH + D \;\approx\; 2DNH + 2DKH \\
\text{MLP:}  \quad & 3DF
\end{align*}

\begin{align*}
\text{Total} &= L(2DNH + 2DKH + 3DF) + DV \\
&= 64 \cdot (2 \cdot 4096 \cdot 32 \cdot 256 + 2 \cdot 4096 \cdot 8 \cdot 256 + 3 \cdot 4096 \cdot 16384) \\
&\quad + 4096 \cdot 32128 \;\approx\; 18.4\text{B params.}
\end{align*}

\subsection*{(ii)}

\[
\text{KV cache size/token} = 2 \cdot H \cdot K \cdot L = 2 \cdot 256 \cdot 8 \cdot 64 = 262{,}144 \text{ bytes} \approx 262 \text{ kB}.
\]

\section*{Q2}

Serve on TPU v5e $4 \times 4$ slice. int8, $T = 128$K.

\subsection*{(i)}

\begin{align*}
\text{KV Cache Size/seq} &= 2HKLT = 2 \cdot 256 \cdot 8 \cdot 64 \cdot 128 \cdot 10^3 \\
&\approx 3.4 \cdot 10^{10} \text{ bytes} = 34 \text{ GB}.
\end{align*}

\[
\text{Model size} = 18.4 \cdot 10^9 \text{ bytes} = 18.4 \text{ GB}.
\]
\[
\text{Total mem} = 16 \cdot 16 \text{ GB} = 256 \text{ GB}.
\]
\[
\text{BS}_{\max} = \frac{256 - 18.4}{34} \approx 7.
\]

\subsection*{(ii)}

Dropping $K = 1$ scales down KV cache size/token by $\frac{1}{8} \Rightarrow \text{BS}_{\max} = 8 \cdot 7 = 56$.

\section*{Q3}

\[
\frac{\text{Model Size}}{(\text{\# devices}) \cdot \Whbm} = \frac{18.4 \cdot 10^9}{16 \cdot 8.2 \cdot 10^{11}} \approx 1 \text{ ms} \quad \text{(assuming full utilization of HBM)}.
\]

\section*{Q4}

What does ICI look like on a $4 \times 4$? $\Rightarrow$ 2 hardware axes, doubled bw.

Roofline bound on TP:
\[
\frac{2 \cdot F}{\left(\frac{C}{\Wici}\right)} = \frac{2 \cdot 16384 \cdot 9.0 \cdot 10^{10}}{3.94 \cdot 10^{14}} = 7 < Y = 16.
\]

We are comm-bounded with TP only.

\paragraph{Prefill:} From our TP roofline, if we use the whole slice for TP, we will be comm-bounded. What I would do is use 8-way TP which will still be slightly comm-bounded but right on the edge, with 2-way sequence parallelism.

\paragraph{Generation:} For generation, our only option is model parallelism. Per layer, since a majority of the parameters are within the MLP (approximating by ignoring einsum):
\[
T_{\text{hbm comm}} = \frac{2DF}{16 \cdot \Whbm}, \qquad T_{\text{ici comm}} = \frac{2BD}{\Wici}.
\]

Want:
\[
T_{\text{ici comm}} < T_{\text{hbm comm}} \;\Rightarrow\; \frac{2BD}{\Wici} < \frac{2DF}{16 \cdot \Whbm}
\]
\[
B < \frac{F \cdot \Wici}{16 \cdot \Whbm} = \frac{16384 \cdot 9.0 \cdot 10^{10}}{16 \cdot 8.2 \cdot 10^{11}}
\]
\[
\Rightarrow B < 112
\]
Which will probably be true in practice since 112 is a pretty big batch size of requests to arrive at once.

For generation, I'd use 16-way TP since we're already comm-bounded and as shown above at reasonable inference batch sizes our ici comm time will be less than the time to load from HBM $\to$ MXU. Since $K$ (\# of KV heads) is $< 16$ but divides it, I'd shard the KV cache along the head and batch dimensions.

Under this sharding each chip would store $\frac{1}{16}$ of all the parameters so assuming BS = 1:
\[
\text{Step Time}_{\min} = \frac{(34 + 18.4) \cdot 10^9}{16 \cdot 8.2 \cdot 10^{11}} \approx 3 \text{ ms}.
\]

This comes from the fact that we're still bottlenecked by loading params from HBM and each device has roughly $(1/16)$ of the data.

\section*{Q5}

$E = 16$, $k = 2$.

\subsection*{(i)}

Per layer:
\begin{align*}
\text{Attn} &= 2DNH + 2DKH \\
\text{MLP}_{\text{total}} &= 3DF \cdot E \\
\text{MLP}_{\text{act}} &= 3DF \cdot k
\end{align*}

\begin{align*}
\text{Total} &= 2DV + L \cdot (2DNH + 2DKH + 3DFE) \\
&= 2 \cdot 4096 \cdot 32128 + 64 \cdot (2 \cdot 4096 \cdot 32 \cdot 256 \\
&\quad + 2 \cdot 4096 \cdot 8 \cdot 256 + 3 \cdot 4096 \cdot 16384 \cdot 16) \\
&\approx 211\text{B}.
\end{align*}
\begin{align*}
\text{Activated} &= 2DV + L \cdot (2DNH + 2DKH + 3DFk) \\
&\approx 31\text{B}.
\end{align*}

\subsection*{(ii)}

Since MLP accounts for most of the FLOPs, we see the bs on a single device where we become comm bound for one of the matmuls (all 3 are the same size). Also assuming we're in generation and our sequence dim is trivial.

\[
\text{AI} = \frac{2BDFk}{2BD + 2DFE + 2BFk} \approx \frac{2BDFk}{2DFE} = B \cdot \frac{k}{E}.
\]

\[
B \cdot \frac{k}{E} \geq \text{Intensity}(\text{v5e}) = 240
\]

\[
\Rightarrow B \geq \frac{240 \cdot E}{k} = 1920.
\]

\subsection*{(iii)}

The KV Cache size per token stays the same as 262 kB since MoE doesn't change the K, V projections.

\subsection*{(iv)}

Again, our linear projections dominate FLOPs so:
\[
\text{FLOPS} \approx 2 \cdot (\text{activated params}) \cdot T = 2 \cdot 31 \cdot 10^9 \cdot T = 62 \cdot 10^9 \cdot T.
\]

\section*{Q6}

Each FFW weight sharded as $[E_z, D_x, F_y]$ ($z$ only used during training as FSDP dim).

\subsection*{(i)}

TPU v5e $8 \times 16$ slice, $Y = 8$, $Z = 16$, int 8.
\[
T_{\text{weights}} \approx \frac{211 \cdot 10^9}{8 \cdot 16 \cdot \Whbm} \approx 2 \text{ ms} \quad \text{(most parameters are in the MLPs so this is a good approximation)}.
\]

\[
\text{Per chip mem used} = \left( \frac{211 \cdot 10^9}{8 \cdot 16} \right) \approx 1.65 \text{ GB}.
\]
\[
\text{Remaining per chip mem} = 16 - 1.65 = 14.35 \text{ GB}.
\]

\subsection*{(ii)}

Again, assuming weights are stored in int8:
\[
N_{\min} = \frac{211}{16} = 14 \text{ v5es} \;\Rightarrow\; 7 \times 2 \text{ slice} \quad \text{(no wraparound connections though)}.
\]

\section*{Q7}

For 2D sharding:
\[
\Tmath = \left( \frac{2BDF}{XYZ} + \frac{2BDF}{XYZ} \right) \bigg/ C = \frac{4BDF}{CXYZ}.
\]

Assuming bf16, bidirectional ici:
\[
\Tcomm = \frac{2B \cdot \left( \frac{D}{X} \right)}{2 \cdot \Wici} + \frac{2 \cdot 2B \cdot \left( \frac{F}{YZ} \right)}{\Wici} + \frac{2B \cdot \left( \frac{D}{X} \right)}{2 \cdot \Wici}
\]
\[
= \frac{4BD}{2X\Wici} + \frac{4BF}{YZ\Wici}.
\]

Since we are free to set the topology, we want to minimize $X, YZ$ (only the product matters for $\Tcomm$):
\[
YZ = N/X, \qquad X = N/YZ.
\]
\[
\Tcomm = \frac{4BD}{2X\Wici} + \frac{4BFX}{N\Wici}.
\]
\[
\frac{\partial \Tcomm}{\partial X} = \frac{-4BD}{2\Wici X^2} + \frac{4BF}{N\Wici} = 0
\]
\[
\Rightarrow \frac{4BF}{N\Wici} = \frac{4BD}{2\Wici X^2} \;\Rightarrow\; X^2 = \frac{D}{2\Wici} \cdot \frac{N\Wici}{F} = \frac{ND}{2F}.
\]
\[
X_{\text{opt}} = \sqrt{\frac{ND}{2F}}, \qquad YZ_{\text{opt}} = \frac{N}{X} = \sqrt{\frac{N^2 \cdot 2F}{ND}} = \sqrt{\frac{2FN}{D}}.
\]

Assuming $F = 4D$ like in specified model:
\[
X_{\text{opt}} = \sqrt{\frac{ND}{2 \cdot 4D}} = \sqrt{\frac{N}{8}}.
\]

\begin{align*}
\Tcomm &= \frac{4BD}{2 \cdot \sqrt{N/8} \cdot \Wici} + \frac{4BF \cdot \sqrt{N/8}}{N \cdot \Wici} \\
&= \frac{2BD}{\Wici \cdot \sqrt{N/8}} + \frac{16BD \cdot \sqrt{N/8}}{N \cdot \Wici} \\
&= \frac{2BD}{\Wici} \cdot \left( \sqrt{\frac{8}{N}} + 8 \cdot \sqrt{\frac{N}{8N^2}} \right) \\
&= \frac{2BD}{\Wici} \left( \sqrt{\frac{8}{N}} + \sqrt{\frac{64N}{8N^2}} \right) = \frac{4BD}{\Wici} \cdot \sqrt{\frac{8}{N}} \\
&= \frac{\sqrt{128} \cdot BD}{\Wici \cdot \sqrt{N}}.
\end{align*}

For normal 3D model parallelism ($[D, F_{XYZ}]$):
\[
\Tmath = \frac{2BDF}{XYZ \cdot C} \cdot 2 = \frac{4BDF}{XYZ \cdot C}.
\]
\[
\Tcomm = \frac{2 \cdot 2 \cdot BD}{3 \cdot \Wici} = \frac{4BD}{3\Wici}.
\]

The time for math comp is the same so we're looking for when:
\[
\Tcomm^{\,3D} > \Tcomm^{\,2D} \;\Rightarrow\; \frac{4BD}{3\Wici} > \frac{\sqrt{128} \cdot BD}{\Wici \cdot \sqrt{N}}
\]
\[
\Rightarrow \sqrt{N} > \frac{3 \cdot \sqrt{128}}{4} \;\Rightarrow\; N > \frac{9 \cdot 128}{16} = 72.
\]

When we have more than 72 chips in our slice (assuming slice is chosen s.t. we have bidirectional ici over all axes), the 2D scheme will outperform the 3D one.

\end{document}
